% !TeX root = main.tex
% =============================================================================
% SEÇÃO 1: CONTEXTO E MOTIVAÇÃO
% =============================================================================
% Esta seção deve apresentar:
% - O contexto geral do projeto de pesquisa
% - A motivação e relevância do tema escolhido
% - O problema de pesquisa identificado
% - Os objetivos gerais e específicos do projeto
% - Referências teóricas que fundamentam a pesquisa
% =============================================================================

\section{Contexto e Motivação}

% --- Introdução ao Tema ---
% Apresente o contexto geral da área de pesquisa e sua relevância.
% Utilize referências bibliográficas para fundamentar suas afirmações.

[Apresente aqui o contexto geral da pesquisa, incluindo a relevância do tema 
e dados que justifiquem a importância do estudo. Utilize referências 
bibliográficas apropriadas para fundamentar suas afirmações.]

% --- EXEMPLOS DE CITAÇÕES ---
% Use \cite{chave} para citações numéricas: [1]
% Exemplo: Segundo a literatura \cite{autor2023_exemplo}, este método é eficaz.
%
% Use \citeonline{chave} quando o autor faz parte da sentença:
% Exemplo: \citeonline{autor2022_livro} demonstrou que...
%
% Para múltiplas referências: \cite{autor2023_exemplo,autor2022_livro}
% Resultado: [1, 2]

% --- Problema de Pesquisa ---
% Descreva o problema específico que sua pesquisa pretende abordar.

[Descreva o problema ou lacuna identificada na área de pesquisa que motivou 
este projeto de iniciação científica.]

% --- Objetivos da Pesquisa ---
% Apresente os objetivos gerais e específicos do projeto.

[Apresente os objetivos do projeto, incluindo o objetivo geral e os 
objetivos específicos que orientam o desenvolvimento da pesquisa. 
Seja claro e objetivo na definição de cada meta a ser alcançada.]

% --- Referencial Teórico (Opcional) ---
% Se necessário, inclua uma breve revisão de literatura ou fundamentação teórica.

[Opcionalmente, apresente uma breve revisão dos principais trabalhos e 
conceitos teóricos que fundamentam sua pesquisa, destacando como eles 
se relacionam com seus objetivos.]

% Exemplo de texto com citações:
% A metodologia proposta por \citeonline{autor2023_exemplo} vem sendo amplamente
% utilizada em diversos contextos \cite{autor2022_livro,autor2024_conf}. Estudos
% recentes \cite{autor2023_tese} demonstram a eficácia dessa abordagem.